% Part: second-order-logic
% Chapter: syntax-and-semantics
% Section: size-of-domain

\documentclass[../../../include/open-logic-section]{subfiles}

\begin{document}

% SEGMENT OLP-0329-S01

\olfileid{sol}{syn}{siz}

\olsection{முடிவிலா மற்றும் \usetoken{S}{enumerable}
  \usetoken{P}{domain} ஐ விவரித்தல்}

!!a{injective} சார்பு~$f\colon M \to M$ ஒன்று !!{surjective} ஆகாமல்,
அதாவது $\ran{f} \neq M$ ஆக இருந்தால், எனில் மட்டுமே, ஒரு கணம்~$M$
(Dedekind) முடிவிலாதது. முதல்தரத் தருக்கத்தில் ஓரிட
!!{function}~$f$ ஐக் கருதி, !!a{structure}~$\Struct{M}$ இல் அதற்கு
அளிக்கப்பட்ட சார்பு~$\Assign{f}{M}$ !!{injective} என்றும்
$\ran{f} \neq \Domain{M}$ என்றும் கூறலாம்:
\[
\lforall[x][\lforall[y][(\eq[f(x)][f(y)] \to \eq[x][y])]] \land
\lexists[y][\lforall[x][\eq/[y][f(x)]]].
\]
$\Struct{M}$ இந்த !!{sentence} ஐ நிறைவுறுத்தினால்,
$\Assign{f}{M}: \Domain{M} \to \Domain{M}$ !!{injective}; எனவே
$\Domain{M}$ முடிவிலாததாக இருக்க வேண்டும். $\Domain{M}$ முடிவிலாதது
எனில் இத்தகைய சார்பு ஒன்று உண்டு; $\Assign{f}{M}$ ஐ அந்தச் சார்பாக
அமைத்தால் $\Struct{M}$ அந்த !!{sentence} ஐ நிறைவுறுத்தும். ஆனால்
இந்த நோக்கத்திற்காகப் பயன்படுத்தப்படும் தருக்கமல்லாத குறி~$f$ மொழியில்
இருக்க வேண்டும். இரண்டாம்தரத் தருக்கத்தில், இத்தகைய சார்பு
\emph{உள்ளது} என்று நேரடியாகக் கூறலாம். இதற்கு இனி~$f$ தேவையில்லை;
தூய இரண்டாம்தரத் தருக்கத்தில் பின்வரும் !!{sentence} கிடைக்கிறது:
\[
\fn{Inf} \ident \lexists[u][(\lforall[x][\lforall[y][(\eq[u(x)][u(y)]
      \lif \eq[x][y])]] \land \lexists[y][\lforall[x][\eq/[y][u(x)]]])].
\]
$\Sat{M}{\fn{Inf}}$ என்பது $\Domain{M}$ முடிவிலாதது எனில், எனில்
மட்டுமே. பின்னர் $\fn{Fin} \ident \lnot \fn{Inf}$ என்று
வரையறுக்கலாம்; $\Sat{M}{\fn{Fin}}$ என்பது $\Domain{M}$ முடிவுறுவது
எனில், எனில் மட்டுமே. தூய முதல்தரத் தருக்கத்தின் எந்த ஒற்றை
!!{sentence} உம் !!{domain} முடிவிலாதது என்பதை வெளிப்படுத்த முடியாது;
ஆனால் அவற்றின் ஒரு முடிவிலா கணம் அதை வெளிப்படுத்த முடியும்.
!!a{structure} இல் அதன் பொருட்களம் முடிவுறுவது அப்போதும் அப்போது
மட்டுமே நிறைவுறுத்தப்படும் தூய முதல்தர !!{sentence}s இன் கணம்
எதுவும் இல்லை.

\begin{prop}
$\Sat{M}{\fn{Inf}}$ என்பது $\Domain{M}$ முடிவிலாதது எனில், எனில் மட்டுமே.
\end{prop}

\begin{proof}
$\Sat{M}{\fn{Inf}}$ என்பது ஏதோ ஓர்~$s$ க்கு
  $\Sat{M}{\lforall[x][\lforall[y][(\eq[u(x)][u(y)] \lif \eq[x][y])]]
    \land \lexists[y][\lforall[x][\eq/[y][u(x)]]]}[s]$ எனில், எனில்
  மட்டுமே. அது பொருந்தினால், $s(u)$ ஓர் !!a{injective} சார்பு;
  ஏதோ ஒரு $y \in \Domain{M}$ ஆனது~$s(u)$ இன் வீச்சில் இல்லை.
  மறுதிசையில், $\ran{f} \neq \Domain{M}$ ஆகும் ஓர் !!a{injective}
  $f\colon \Domain{M} \to \Domain{M}$ இருந்தால், $s(u) = f$ என்பது
  தேவையான மாறி மதிப்பீடு.
\end{proof}

ஒரு கணம் $M$ இன் !!{element}s க்கு (மறுநிகழ்வு இல்லாத, ஆனால்
முடிவுறக்கூடிய) பின்வரும் ஓர் எண்ணிடல் இருந்தால், அந்தக் கணம்
!!{enumerable}:
\[
m_0, m_1, m_2, \dots
\]
இத்தகைய எண்ணிடல் இருப்பது அப்போதும் அப்போது மட்டுமே,
!!a{element} $z \in M$ ஒன்றும் $f\colon M \to M$ என்ற சார்பும் இருந்து,
$z$, $f(z)$, $f(f(z))$, \dots ஆகியவை~$M$ இன் எல்லா
!!{element}s ஆகும். எண்ணிடல் இருந்தால், $z = m_0$ என்றும்
$f(m_k) = m_{k+1}$ என்றும் (அல்லது $m_k$ எண்ணிடலின் கடைசி
!!{element} எனில் $f(m_k) = m_k$ என்றும்) அமைத்தால், தேவையான
!!{element} உம் சார்பும் கிடைக்கும். மறுபுறம், இத்தகைய $z$, $f$
இருந்தால், $z$, $f(z)$, $f(f(z))$, \dots என்பது~$M$ இன் எண்ணிடல்;
எனவே $M$ !!{enumerable}. இரண்டாம்தரத் தருக்கத்தில் $z$,~$f$
இருப்பதை வெளிப்படுத்தி, ஒரு !!a{structure} இல் அந்த !!{structure}
!!{enumerable} ஆக இருந்தால், எனில் மட்டுமே, மெய்யாகும் பின்வரும்
!!{sentence} ஐ உருவாக்கலாம்:
\[
\fn{Count} \ident
\lexists[z][\lexists[u][\lforall[X][((X(z) \land
      \lforall[x][(X(x) \lif X(u(x)))]) \lif \lforall[x][X(x)])]]]
\]

\begin{prop}
$\Sat{M}{\fn{Count}}$ என்பது $\Domain{M}$ !!{enumerable} எனில்,
எனில் மட்டுமே.
\end{prop}

\begin{proof}
$\Domain{M}$ !!{enumerable} என்று கருதி, $m_0$, $m_1$, \dots
என்பது ஓர் எண்ணிடல் என்க. மறுநிகழ்வுகளை அகற்றி எந்த $m_k$ உம்
இருமுறை வராது என்று உறுதிசெய்யலாம். $f(m_k) = m_{k+1}$ என்று
வரையறுத்து, $s(z) = m_0$, $s(u) = f$ எனக் கொள்க. பின்வருவதை
நாம் காட்டுகிறோம்:
\[
\Sat{M}{\lforall[X][((X(z) \land \lforall[x][(X(x) \lif X(u(x)))])
    \lif \lforall[x][X(x)])]}[s]
\]
$M \subseteq \Domain{M}$ தன்னிச்சையானது என்க. மேலும்
$\Sat{M}{(X(z) \land \lforall[x][(X(x) \lif
X(u(x)))])}[\Subst{s}{M}{X}]$ என்று கருதுக. அப்போது
$\Subst{s}{M}{X}(z) \in M$; மேலும் $x \in M$ ஆகும்போதெல்லாம்
$(\Subst{s}{M}{X}(u))(x) \in M$. வேறு சொற்களில்,
$\varAssign{\Subst{s}{M}{X}}{s}{X}$ என்பதால், $m_0 \in M$;
மேலும் $x \in M$ எனில் $f(x) \in M$. எனவே $m_0 \in M$,
$m_1 = f(m_0) \in M$, $m_2 = f(f(m_0)) \in M$ முதலியன.
ஆகவே $M = \Domain{M}$; எனவே
$\Sat{M}{\lforall[x][X(x)]}[\Subst{s}{M}{X}]$. மேலும்
$M \subseteq \Domain{M}$ தன்னிச்சையானது என்பதால், தேவையான
$\Sat{M}{\fn{Count}}$ கிடைக்கிறது.

இப்போது $\Sat{M}{\fn{Count}}$ என்று, அதாவது ஏதோ ஓர்~$s$ க்கு
\[
\Sat{M}{\lforall[X][((X(z) \land \lforall[x][(X(x) \lif X(u(x)))])
    \lif \lforall[x][X(x)])]}[s]
\]
என்று கருதுக. $m = s(z)$, $f = s(u)$ எனக் கொடுத்து,
$M = \{m, f(m), f(f(m)), \dots\}$ ஐக் கருதுக. இவ்வாறு
வரையறுக்கப்பட்ட $M$ தெளிவாக !!{enumerable}. அனுமானத்தின்படி,
\[
\Sat{M}{(X(z) \land \lforall[x][(X(x) \lif X(u(x)))])
    \lif \lforall[x][X(x)]}[\Subst{s}{M}{X}]
\]
மேலும், $M \ni m = \Subst{s}{M}{X}(z)$ என்பதால்
$\Sat{M}{X(z)}[\Subst{s}{M}{X}]$; $x \in M$ ஆகும்போதெல்லாம்
$f(x) \in M$ என்பதால்
$\Sat{M}{\lforall[x][(X(x) \lif
X(u(x)))]}[\Subst{s}{M}{X}]$ உம் பொருந்தும். முன்னிகழ்வும்
நிபந்தனைவாக்கியமும் இரண்டும் நிறைவுறுத்தப்படுவதால், பின்னிகழ்வும்
நிறைவுறுத்தப்பட வேண்டும்:
$\Sat{M}{\lforall[x][X(x)]}[\Subst{s}{M}{X}]$. ஆனால் இதன் பொருள்
$M = \Domain{M}$. எனவே, வரையறைப்படி $M$ எண்ணிடத்தக்கது என்பதால்
$\Domain{M}$ உம் !!{enumerable}.
\end{proof}

\begin{prob}
!!{sentence}~$\fn{Inf} \land \fn{Count}$ ஆனது எல்லா
!!{denumerable} பொருட்களங்களிலும், அவற்றில் மட்டுமே, மெய்யாகும்.
~$\fn{Count}$ இன் வரையறையை மாற்றி, பொருட்களம் !!{denumerable}
என்பதை நேரடியாக வெளிப்படுத்தும் வேறொரு !!{sentence} ஆக அதை
அமைத்து, அது அவ்வாறு வெளிப்படுத்துகிறது என்று நிறுவுக.
\end{prob}

\end{document}
